\section{The planar rigidity matroid}
\label{sec:rigidity-matroid}

\Cref{sec:laman-theorem} proved the ``easy direction'' of
Lov\'asz--Yemini~\cite{lovaszYemini1982}: in dimension $2$, a
row-independent (\cref{def:edgeSet-rowIndependent}) edge set is
automatically $(2, 3)$-sparse
(\cref{lem:isSparse-of-rowIndependent-two}). This chapter completes
the matroid identification by proving the harder converse --- every
$(2, 3)$-sparse edge set is row-independent at some placement --- and
packages the resulting independence structure as a matroid on the
off-diagonal symmetric pairs over $V$.

The proof follows the classical strategy of Lov\'asz--Yemini, in the
modern presentation of Jord\'an~\cite{jordan2016}~\S 2.2 (Theorem
2.2.1): induction on $|E|$, with each Henneberg move on a sparse graph
lifting row-independence to a fresh placement of the larger graph by
placing the new vertex generically. The combinatorial reverse step
(Jord\'an Lemma 2.1.4) --- a sparse graph on $|V| \ge 1$ with $|E| \ge 1$
admits either a Type~I reverse (delete a degree-${\le}2$ vertex) or a
Type~II reverse (split a degree-$3$ vertex) yielding a smaller sparse
graph --- is the sparse-graph analogue of the Laman reverse decomposition
(\cref{thm:isLaman-exists-typeI-or-typeII-reverse}) and reuses the same
tight-subset machinery (\cref{thm:isTightOn-union-inter}).

\emph{Terminology.} Throughout this chapter we use the project's
\emph{$(2, 3)$-tight subset} (\cref{def:isTightOn}) in place of
Jord\'an's \emph{critical set}. The two are interchangeable: a
critical set in a $(2, 3)$-sparse graph $G$ is exactly a subset
$X \subseteq V$ with $\bigl|\edgesIn{X}\bigr| = 2|X| - 3$. Likewise
Jord\'an's Lemma~2.1.2
(critical sets are union/intersection-closed when intersections have
$\ge 2$ vertices) is our \cref{thm:isTightOn-union-inter}; the
three-pair tight-union (Jord\'an Lemma~2.1.3) is
\cref{thm:isTightOn-union-with-bonus} composed twice.

\subsection{Sparse-graph reverse decomposition}
\label{sec:rigidity-matroid-sparse-reverse}

\begin{theorem}[Sparse-graph reverse decomposition]
  \label{thm:isSparse-exists-typeI-or-typeII-reverse}
  \lean{SimpleGraph.IsSparse.exists_typeI_or_typeII_reverse}
  \leanok
  \uses{def:isSparse, def:typeI, def:typeII}
  Let $G$ be a $(2, 3)$-sparse graph on a finite vertex set $V$ with at
  least one edge. Then there exists $v \in V$ with
  $d_{G}(v) \in \{1, 2, 3\}$ admitting one of three reverses:
  \begin{enumerate}
    \item ($d_{G}(v) = 1$, \emph{pendant reverse}) $v$ has a unique
      neighbour $a$, and $G - v$ is $(2, 3)$-sparse.
    \item ($d_{G}(v) = 2$, \emph{Type~I reverse}) $v$ has two distinct
      neighbours $a \ne b$, and $G - v$ is $(2, 3)$-sparse.
    \item ($d_{G}(v) = 3$, \emph{Type~II reverse}) $v$ has three
      neighbours $\{u, w, z\} = N_{G}(v)$ with a non-adjacent pair
      (say $\{u, w\}$) such that the graph $G^{u,w}_{v}$ obtained by
      deleting $v$ and adding the edge $\{u, w\}$ is $(2, 3)$-sparse.
  \end{enumerate}
\end{theorem}
\begin{proof}
  \leanok
  \uses{def:isSparse, def:typeI, def:typeII, thm:isTightOn-union-inter,
    lem:isSparse-typeII-reverse-blocker,
    lem:isSparse-exists-nonadj-among-three-neighbors}
  Restrict attention to $V_{+} := \{v \in V : d_{G}(v) \ge 1\}$; the
  endpoints of any edge lie in $V_{+}$, so $|V_{+}| \ge 2$ and
  $E(G[V_{+}]) = E(G)$. Since $G[V_{+}]$ is sparse (subgraph of
  sparse) with $|V_{+}| \ge 2$ vertices and the same edge set, the
  handshake / sparsity bound on $G[V_{+}]$ supplies a vertex $v$ with
  $d_{G[V_{+}]}(v) \le 3$, hence the same $d_{G}(v) \le 3$ (no edges
  to isolated vertices). In particular $d_{G}(v) \in \{1, 2, 3\}$.

  The case $d_{G}(v) \le 2$ yields a pendant or Type~I reverse:
  deleting $v$ removes $d_{G}(v)$ edges and sparsity is preserved
  (sparse minus edges is sparse); the neighbour data is read off from
  $|N_{G}(v)| = d_{G}(v)$ by cardinality.

  The case $d_{G}(v) = 3$
  requires a suitable splitting on a non-adjacent pair of neighbours
  (one exists by \cref{lem:isSparse-exists-nonadj-among-three-neighbors};
  Jord\'an Lemma 2.1.4(b)): if no splitting at $v$ is suitable then the
  per-pair tight-blocker (\cref{lem:isSparse-typeII-reverse-blocker})
  yields three maximal
  $(2, 3)$-tight sets $X_{uw}, X_{uz}, X_{wz} \subseteq V \setminus
  \{v\}$ each containing exactly two neighbours of $v$, with pairwise
  singleton intersections; the tight-union closure
  (\cref{thm:isTightOn-union-inter}) and the three-pair tight-union
  analogue (Jord\'an Lemma 2.1.3) make $X_{uw} \cup X_{uz} \cup
  X_{wz}$ tight, forcing $d_{G}(v, X_{uw} \cup X_{uz} \cup X_{wz})
  \ge 3$, which then contradicts the sparsity bound on $X_{uw} \cup
  X_{uz} \cup X_{wz} \cup \{v\}$.
\end{proof}

\begin{fmlnote}
  \label{fmlnote:reverse-vs-lift-form}
  The reverse decomposition above and the forward lifts below are
  stated in complementary forms. The reverse is \emph{flat}: it
  describes the smaller graph directly by its edges (``$G - v$ is
  sparse'', ``$G^{u,w}_v$ is sparse''), without naming a Henneberg
  constructor, because it takes an arbitrary $G$ and produces a
  smaller witness --- the named operation is incidental. Each forward
  lift is instead a statement \emph{about} a named extension: that
  $\mathrm{typeI}\,G'\,a\,b$ or $\mathrm{typeII}\,G'\,a\,b\,c$
  preserves row-independence, so the operation belongs in the
  statement. The same split governs the rigidity-preservation lemmas
  (\cref{thm:typeI-isGenericallyRigidInj-two},
  \cref{thm:typeII-isGenericallyRigidInj-two}), stated in operation
  form, and the Laman reverse decomposition
  (\cref{thm:isLaman-exists-typeI-or-typeII-reverse}), stated flat.
  The two forms meet at callsites combining a reverse decomposition
  with a forward lift: the flat reverse produces a
  vertex-with-neighbourhood witness, from which the corresponding
  $\mathrm{typeI}$ / $\mathrm{typeII}$ extension and its isomorphism
  to $G$ are reconstructed by the project's isomorphism constructors
  (\texttt{typeI\_iso\_of\_two\_neighbors},
  \texttt{typeII\_iso\_of\_three\_neighbors}).
\end{fmlnote}

\subsection{Lifting row-independence through Henneberg moves}
\label{sec:rigidity-matroid-lifts}

The two Henneberg lifts are the inductive step of the converse. Each
states: given a placement $p'$ of a smaller graph $G'$ at which the
edge set is row-independent, the larger graph $G$ obtained by a
Henneberg move admits a placement at which its edge set is
row-independent.

\begin{lemma}[Type~I row-independence: conditional core]
  \label{lem:typeI-rowIndependent-extend}
  \lean{SimpleGraph.Henneberg.typeI_edgeSetRowIndependent_extend}
  \leanok
  \uses{def:typeI, def:edgeSet-rowIndependent}
  Let $G'$ be a graph on a finite vertex set $V$ and $p' : V \to \R^{2}$
  a placement at which $E(G')$ is row-independent. Let $a, b \in V$ and
  $q \in \R^{2}$, and suppose $q - p'_a$ and $q - p'_b$ are linearly
  independent. Then the extended placement
  $\hat p : \mathrm{Option}\,V \to \R^{2}$ with
  $\hat p(\mathrm{none}) = q$ and $\hat p(\mathrm{some}\,u) = p'_u$ makes
  $E(\mathrm{typeI}\,G'\,a\,b)$ row-independent.
\end{lemma}
\begin{proof}
  \leanok
  Row-independence analogue of
  \cref{thm:typeI-isInfinitesimallyRigid-extend}. Partition
  $E(\mathrm{typeI}\,G'\,a\,b)$ into the image of $E(G')$ under
  $\mathrm{some}$ (``old'' edges) and the two new edges
  $\{\mathrm{none}, \mathrm{some}\,a\}$, $\{\mathrm{none}, \mathrm{some}\,b\}$.
  The old rigidity-rows at $\hat p$ are pullbacks under the restriction
  $(\cdot \circ \mathrm{some}) : \mathrm{Framework}\,(\mathrm{Option}\,V)\,2
  \to \mathrm{Framework}\,V\,2$ of the $G'$ rigidity-rows at $p'$;
  surjectivity of this restriction makes its dual injective, so the
  row-independence of the $G'$ family transfers.

  The remaining linear-independence and disjointness facts both follow
  from a parametric \emph{test motion} $x_\alpha : \mathrm{Option}\,V
  \to \R^{2}$ defined by $x_\alpha(\mathrm{none}) = \alpha$ and
  $x_\alpha \circ \mathrm{some} = 0$. Evaluating any new row at
  $x_\alpha$ produces $\langle q - p'_a, \alpha \rangle$ or
  $\langle q - p'_b, \alpha \rangle$; old rows evaluate to zero (their
  inner-product arguments only see $x_\alpha \circ \mathrm{some} = 0$),
  and the latter property extends to the entire old-row span by
  linearity. For linear independence of the new rows, a vanishing
  combination $c \cdot \mathrm{row}_a + d \cdot \mathrm{row}_b$ evaluated
  at $x_\alpha$ reads $c \langle q - p'_a, \alpha\rangle + d \langle q -
  p'_b, \alpha \rangle = 0$ for all $\alpha$, forcing $c(q - p'_a) + d(q -
  p'_b) = 0$ and hence $c = d = 0$ by hypothesis. For disjointness of the
  spans, a joint element with new-row representation $c \cdot
  \mathrm{row}_a + d \cdot \mathrm{row}_b$ vanishes on every $x_\alpha$
  (because it lies in the old-row span), so the same
  coefficient-extraction forces $c = d = 0$ and the element is zero.
\end{proof}

\begin{lemma}[Type~I lift of row-independence]
  \label{lem:typeI-rowIndependent-lift}
  \lean{SimpleGraph.Henneberg.typeI_edgeSetRowIndependent_lift}
  \leanok
  \uses{def:typeI, def:edgeSet-rowIndependent,
    lem:typeI-rowIndependent-extend}
  Let $G'$ be a graph on a finite vertex set $V$ and $p' : V \to \R^{2}$
  a placement at which $E(G')$ is row-independent. Let $a, b \in V$ be
  two old vertices with $p'_a \ne p'_b$. Then there exists a placement
  $p : \mathrm{Option}\,V \to \R^{2}$ with $p \circ \mathrm{some} = p'$
  such that $E(\mathrm{typeI}\,G'\,a\,b)$ is row-independent at $p$.
\end{lemma}
\begin{proof}
  \leanok
  \uses{lem:typeI-rowIndependent-extend}
  Pick $q$ off both the lines through $p'_a$ and through $p'_b$ in the
  direction $p'_b - p'_a$ (equivalently, $q - p'_a$ and $q - p'_b$
  linearly independent); the bad set is a finite union of one or two
  lines in $\R^{2}$, so a generic $q$ works.
  \Cref{lem:typeI-rowIndependent-extend} applies and gives
  row-independence of $E(\mathrm{typeI}\,G'\,a\,b)$ at the extended
  placement.
\end{proof}

The pendant reverse of a degree-$1$ vertex is handled by the
degenerate Type~I move $\mathrm{typeI}\,G'\,a\,a$ ($b = a$), which
joins the new vertex to the single old vertex $a$ --- the pendant
attachment used by Lov\'asz--Yemini.

\begin{lemma}[Pendant lift of row-independence]
  \label{lem:pendant-rowIndependent-lift}
  \lean{SimpleGraph.Henneberg.typeI_pendant_edgeSetRowIndependent_lift}
  \leanok
  \uses{def:edgeSet-rowIndependent}
  Let $G'$ be a graph on a finite vertex set $V$ and $p' : V \to
  \R^{2}$ a placement at which $E(G')$ is row-independent. Let
  $a \in V$ be an old vertex. Then there exists a placement
  $p : \mathrm{Option}\,V \to \R^{2}$ with $p \circ \mathrm{some} = p'$
  such that $E(\mathrm{typeI}\,G'\,a\,a)$ is row-independent at $p$.
\end{lemma}
\begin{proof}
  \leanok
  \uses{def:edgeSet-rowIndependent}
  Pick $q \in \R^{2}$ with $q \ne p'_{a}$ (such $q$ exists since
  $\R^{2}$ is uncountable). The edge set of $\mathrm{typeI}\,G'\,a\,a$
  partitions into the image of $E(G')$ under $\mathrm{some}$ and the
  single new edge $\{\mathrm{none}, \mathrm{some}\,a\}$. The old
  rigidity-rows at the extended $\hat p$ pull back under the restriction
  $(\cdot \circ \mathrm{some})$, whose dual map is injective (the
  restriction is surjective), so the old rows' row-independence
  transfers. The single new row is non-zero (witnessed by the motion
  $\mathrm{none} \mapsto q - p'_{a}, \mathrm{some}\,u \mapsto 0$, which
  evaluates the new row to $\|q - p'_{a}\|^2 > 0$). Disjointness of the
  spans uses the same test motion $x_\alpha(\mathrm{none}) = \alpha,
  x_\alpha \circ \mathrm{some} = 0$ as in
  \cref{lem:typeI-rowIndependent-extend}: any joint element vanishes
  at $x_\alpha$, forcing $c \cdot \langle q - p'_{a}, \alpha \rangle
  = 0$ for all $\alpha$, hence $c = 0$.
\end{proof}

\begin{lemma}[Type~II row-independence: conditional core]
  \label{lem:typeII-rowIndependent-extend}
  \lean{SimpleGraph.Henneberg.typeII_edgeSetRowIndependent_extend}
  \leanok
  \uses{def:typeII, def:edgeSet-rowIndependent}
  Let $G'$ be a graph on a finite vertex set $V$ with
  $\{a, b\} \in E(G')$, and $p' : V \to \R^{2}$ a placement at which
  $E(G')$ is row-independent. Let $a, b, c \in V$ and
  $q \in \R^{2}$ with $q$ collinear with $p'_a, p'_b$ but distinct
  from both --- that is, $q - p'_a = s \cdot (p'_b - p'_a)$ for some
  $s \ne 0, 1$ --- and suppose $q - p'_a, q - p'_c$ are linearly
  independent. Then the extended placement
  $\hat p : \mathrm{Option}\,V \to \R^{2}$ with $\hat p(\mathrm{none})
  = q$ and $\hat p(\mathrm{some}\,u) = p'_u$ makes
  $E(\mathrm{typeII}\,G'\,a\,b\,c)$ row-independent.
\end{lemma}
\begin{proof}
  \leanok
  Row-independence analogue of
  \cref{thm:typeII-isInfinitesimallyRigid-extend}. The edge set of
  $\mathrm{typeII}\,G'\,a\,b\,c$ partitions into the image of
  $E(G') \setminus \{\{a, b\}\}$ under $\mathrm{some}$ (``old'' edges)
  and the three new edges $\{\mathrm{none}, \mathrm{some}\,a\}$,
  $\{\mathrm{none}, \mathrm{some}\,b\}$,
  $\{\mathrm{none}, \mathrm{some}\,c\}$. The two new rows for
  $\{\mathrm{none}, \mathrm{some}\,a\}$ and
  $\{\mathrm{none}, \mathrm{some}\,b\}$ are scalar multiples of
  $p'_b - p'_a$ in their $\mathrm{some}\,a$- and
  $\mathrm{some}\,b$-columns: collinearity of $q$ with $p'_a, p'_b$
  forces $q - p'_a = s(p'_b - p'_a)$ and $q - p'_b = (s-1)(p'_b -
  p'_a)$. A direct computation gives the identity
  \[ s \cdot \mathrm{row}_b - (s-1) \cdot \mathrm{row}_a
      = -s(s-1) \cdot T(\mathrm{row}_{G'}\{a, b\}), \]
  where $T$ is the dual (pullback) of the restriction
  $\cdot \circ \mathrm{some}$ (injective since $\mathrm{some}$ is, so
  its dual map is injective via surjectivity of the restriction).
  Thus $\mathrm{span}\{\mathrm{row}_a, \mathrm{row}_b\}$ contains
  $T(\mathrm{row}_{G'}\{a, b\})$, recovering the deleted $G'$-row.

  Linear independence of the three new rows uses two test motions:
  the $\mathrm{none}$-only motion $x_\alpha$ pins
  $sc_a + (s-1)c_b = 0$ and $c_c = 0$; a motion with
  $\mathrm{some}\,a \mapsto q - p'_a$ and all other coordinates zero
  pins $c_a = 0$ (using $s \ne 0$ and $q \ne p'_a$).

  Disjointness between the lifted-old span (the image of $T$ on the
  rows of $E(G') \setminus \{\{a, b\}\}$) and the new-row span runs the
  same test-motion argument to produce $sc_a + (s-1)c_b = 0$ and
  $c_c = 0$; the surviving combination $c_a \cdot \mathrm{row}_a +
  c_b \cdot \mathrm{row}_b$ then equals $s c_a \cdot
  T(\mathrm{row}_{G'}\{a, b\})$ by the identity above. If
  $s c_a \ne 0$, dividing recovers $T(\mathrm{row}_{G'}\{a, b\})$ from
  the lifted-old span, hence (by injectivity of $T$)
  $\mathrm{row}_{G'}\{a, b\}$ from the $G'$-row span on
  $E(G') \setminus \{\{a, b\}\}$, contradicting the row-independence of
  $E(G')$. So $c_a = 0$, then $c_b = 0$ by the first equation, and the
  entire combination vanishes.
\end{proof}

\begin{lemma}[Row-independence is open in the placement]
  \label{lem:edgeSet-rowIndependent-eventually}
  \lean{SimpleGraph.EdgeSetRowIndependent.eventually}
  \leanok
  \uses{def:edgeSet-rowIndependent, def:rigidityRow,
    lem:isInfinitesimallyRigid-eventually}
  Let $G$ be a graph on a finite vertex set $V$. If
  $p_0 : V \to \R^{d}$ makes an edge subset $I \subseteq E(G)$
  row-independent, then so does every placement in some neighborhood
  of $p_0$.
\end{lemma}
\begin{proof}
  \leanok
  Row-independence analogue of
  \cref{lem:isInfinitesimallyRigid-eventually}. Since
  $\mathrm{Module.Dual}\,\R\,(\mathrm{Framework}\,V\,d)$ carries no
  canonical norm, transport the linear-independence assertion to
  $\R^{n}$ (with $n = \dim\,\mathrm{Framework}\,V\,d$) along the basis
  isomorphism $\psi : \mathrm{Dual} \to \R^{n}$, $\psi\,l\,i = l(b\,i)$
  for a fixed basis $b$ of $\mathrm{Framework}\,V\,d$.

  The matrix family
  $p \mapsto (i, k) \mapsto \mathrm{rigidityRow}\,G\,p\,i\,(b\,k)$ is
  continuous in $p$ (each entry is an inner product, jointly continuous
  in $(p, b\,k)$ with $b\,k$ fixed) and equals the $\psi$-image of a
  linearly independent family at $p_0$; linear independence is open in a
  finite-dimensional normed space, so it persists on a neighborhood, and
  transporting back along $\psi^{-1}$ (a linear equivalence preserves
  linear independence in both directions) gives the claim.
\end{proof}

\begin{lemma}[Row-independence transports along a graph isomorphism]
  \label{lem:edgeSet-rowIndependent-iso}
  \lean{SimpleGraph.EdgeSetRowIndependent.iso}
  \leanok
  \uses{def:edgeSet-rowIndependent, def:rigidityRow}
  Let $\varphi : G \simeq H$ be a graph isomorphism between finite
  graphs and $q$ a framework of $H$. If every edge of $H$ is
  row-independent at $q$, then every edge of $G$ is row-independent at
  the pulled-back placement $q \circ \varphi$.
\end{lemma}
\begin{proof}
  \leanok
  Row-independence analogue of \cref{lem:isInfinitesimallyRigid-iso}.
  Precomposition with $\varphi^{-1}$ is an $\R$-linear equivalence
  $L_\varphi : \mathrm{Framework}\,V\,d \xrightarrow{\sim}
  \mathrm{Framework}\,W\,d$, $m \mapsto m \circ \varphi^{-1}$. A direct
  computation on $e = \{u, w\} \in E(G)$ gives
  $\mathrm{rigidityRow}\,G\,(q \circ \varphi)\,e =
  L_\varphi^{*}\bigl(\mathrm{rigidityRow}\,H\,q\,(\varphi\,e)\bigr)$,
  where $\varphi$ acts on edges via $\mathrm{Sym2.map}$. The $H$-row
  family is linearly independent by hypothesis; reindexing along the
  edge-set bijection $\varphi_E : E(G) \xrightarrow{\sim} E(H)$ and
  applying the dual of a linear equivalence (an injective linear map)
  both preserve linear independence.
\end{proof}

\begin{lemma}[Type~II lift of row-independence]
  \label{lem:typeII-rowIndependent-lift}
  \lean{SimpleGraph.Henneberg.typeII_edgeSetRowIndependent_lift}
  \leanok
  \uses{def:typeII, def:edgeSet-rowIndependent,
    lem:typeII-rowIndependent-extend}
  Let $G'$ be a graph on a finite vertex set $V$, $\{u, w\} \in E(G')$,
  and $z \in V$ a third vertex distinct from $u$ and $w$. Let $p' :
  V \to \R^{2}$ be a placement at which $E(G')$ is row-independent.
  Let $G$ be the Type~II extension of $G'$ deleting $\{u, w\}$ and
  adding a fresh vertex $\star$ adjacent to $u, w, z$. Then there
  exists a placement $p$ of $V \sqcup \{\star\}$ at which $E(G)$ is
  row-independent.
\end{lemma}
\begin{proof}
  \leanok
  \uses{lem:typeII-rowIndependent-extend,
    lem:edgeSet-rowIndependent-eventually}
  Apply \cref{lem:typeII-rowIndependent-extend} to a placement $p'$
  perturbed to make $p'(u), p'(w), p'(z)$ non-collinear (otherwise
  the third new row $\{\star, z\}$ would collapse into the span of
  the first two). The perturbation translates $p'(z)$ along a
  direction outside the line through $p'(u), p'(w)$;
  \cref{lem:edgeSet-rowIndependent-eventually} permits the
  perturbation magnitude to be taken arbitrarily small, preserving
  row-independence of $E(G')$.

  With the non-collinear placement in hand, set
  $s := 1/2$ and place $p(\star) := p(u) + s \cdot (p(w) - p(u))$.
  Then $s \ne 0, 1$ and the auxiliary linear independence
  $p(\star) - p(u), p(\star) - p(z)$ follows from non-collinearity
  (row operations on $(s \cdot (p(w) - p(u)), s \cdot (p(w) - p(u))
  - (p(z) - p(u)))$ produce $(p(w) - p(u), p(z) - p(u))$).
  \Cref{lem:typeII-rowIndependent-extend} then applies. The
  $p'(u) \ne p'(w)$ requirement of the perpendicular direction is
  supplied by the row-independence of $E(G')$: a zero rigidity row at
  the edge $\{u, w\}$ would contradict linear independence.
\end{proof}

\subsection{The hard direction}
\label{sec:rigidity-matroid-hard-direction}

\begin{theorem}[$(2, 3)$-sparse $\Rightarrow$ row-independent at some placement]
  \label{thm:isSparse-exists-rowIndependent-placement}
  \lean{SimpleGraph.IsSparse.exists_rowIndependent_placement}
  \leanok
  \uses{def:isSparse, def:edgeSet-rowIndependent,
    thm:isSparse-exists-typeI-or-typeII-reverse,
    lem:pendant-rowIndependent-lift,
    lem:typeI-rowIndependent-lift,
    lem:typeII-rowIndependent-lift,
    lem:edgeSet-rowIndependent-iso}
  Let $H$ be a $(2, 3)$-sparse graph on a finite vertex set $V$. Then
  there exists a placement $p : V \to \R^{2}$ at which $E(H)$ is
  row-independent.
\end{theorem}
\begin{proof}
  \leanok
  \uses{thm:isSparse-exists-typeI-or-typeII-reverse,
    lem:pendant-rowIndependent-lift,
    lem:typeI-rowIndependent-lift,
    lem:typeII-rowIndependent-lift,
    lem:edgeSet-rowIndependent-iso,
    lem:edgeSet-rowIndependent-eventually}
  Strong induction on $|V|$ (each Henneberg reverse strictly decreases
  $|V|$, and the base case below covers $|E(H)| = 0$ uniformly). In the
  base case, when $E(H)$ is empty (in particular when $|V| \le 1$), the
  empty family of rows is linearly independent at any placement.

  For the inductive step,
  \cref{thm:isSparse-exists-typeI-or-typeII-reverse} supplies a
  vertex $v$ with $d_H(v) \in \{1, 2, 3\}$ and a smaller sparse graph
  $H'$ on $V \setminus \{v\}$ via one of three reverses (pendant,
  Type~I, or Type~II); the induction hypothesis gives a row-independent
  placement $p'$ of $H'$. The corresponding lift
  (\cref{lem:pendant-rowIndependent-lift},
  \cref{lem:typeI-rowIndependent-lift}, or
  \cref{lem:typeII-rowIndependent-lift}) reconstructs a row-independent
  placement of the Henneberg extension $\mathrm{typeI}\,H'\,a\,a$,
  $\mathrm{typeI}\,H'\,a\,b$, or $\mathrm{typeII}\,H'\,x\,y\,c$, which
  is graph-isomorphic to $H$ via
  \texttt{typeI\_iso\_of\_two\_neighbors} /
  \texttt{typeII\_iso\_of\_three\_neighbors} (project-internal);
  \cref{lem:edgeSet-rowIndependent-iso} transports row-independence
  back to $H$. The Type~I branch additionally applies a small
  perturbation to $p'$ (combining a step-along-$\delta a$ deformation
  with \cref{lem:edgeSet-rowIndependent-eventually}) to ensure
  $p'(a) \ne p'(b)$, the linear-independence input of
  \cref{lem:typeI-rowIndependent-lift}.
\end{proof}

\begin{theorem}[Lov\'asz--Yemini: row-independence iff sparse, dimension $2$]
  \label{thm:edgeSet-rowIndependent-iff-isSparse}
  \lean{SimpleGraph.edgeSet_rowIndependent_iff_isSparse_dim_two}
  \leanok
  \uses{def:edgeSet-rowIndependent, def:isSparse,
    thm:isSparse-exists-rowIndependent-placement,
    lem:isSparse-of-rowIndependent-two,
    lem:exists-affinelySpanning-of-eventually,
    lem:edgeSet-rowIndependent-eventually}
  Let $G$ be a graph on a finite vertex set $V$. An edge set
  $I \subseteq E(G)$ is row-independent at some placement
  $p : V \to \R^{2}$ if and only if the spanning subgraph of $G$ with
  edge set $I$ is $(2, 3)$-sparse.
\end{theorem}
\begin{proof}
  \leanok
  \uses{thm:isSparse-exists-rowIndependent-placement,
    lem:isSparse-of-rowIndependent-two,
    lem:exists-affinelySpanning-of-eventually,
    lem:edgeSet-rowIndependent-eventually}
  $(\Rightarrow)$ is the easy direction
  \cref{lem:isSparse-of-rowIndependent-two}, after perturbing the
  witness placement to one that affinely spans every triple via
  \cref{lem:exists-affinelySpanning-of-eventually} specialised at
  row-independence (using \cref{lem:edgeSet-rowIndependent-eventually}
  as the openness witness).

  $(\Leftarrow)$ is
  \cref{thm:isSparse-exists-rowIndependent-placement} applied to the
  spanning subgraph $H = \mathrm{fromEdgeSet}(I)$, bridged back to a
  row-independent placement of $E(G)$ on $I$ via the per-edge equality
  of rigidity rows: both reduce to $\langle p u - p v, \mathrm{motion}\,u
  - \mathrm{motion}\,v\rangle$ on the same $\mathrm{Sym}_{2}$-value, so
  the equality is definitional.
\end{proof}

\begin{corollary}[Every Laman graph admits a row-independent placement]
  \label{cor:isLaman-exists-rowIndependent}
  \lean{SimpleGraph.IsLaman.exists_rowIndependent_placement}
  \leanok
  \uses{def:isLaman, def:edgeSet-rowIndependent,
    thm:isSparse-exists-rowIndependent-placement}
  Let $H$ be a Laman graph on a finite vertex set $V$. Then there
  exists a placement $p : V \to \R^{2}$ at which $E(H)$ is
  row-independent.
\end{corollary}
\begin{proof}
  \leanok
  \uses{thm:isSparse-exists-rowIndependent-placement}
  Laman implies $(2, 3)$-sparse, and
  \cref{thm:isSparse-exists-rowIndependent-placement} applies directly.
\end{proof}
\begin{fmlnote}
  The Lean witness specialises
  \cref{thm:isSparse-exists-rowIndependent-placement} in one line; it is
  marked deprecated in favour of that general sparse-level form and
  exists only as the formal anchor for this corollary.
\end{fmlnote}

\subsection{The rigidity matroid (dimension $2$)}
\label{sec:rigidity-matroid-instance}

We are now in a position to define the planar rigidity matroid and
identify its independent sets with the $(2, 3)$-sparse edge subsets
of the complete graph $K_{V}$ on $V$
(Lov\'asz--Yemini~\cite{lovaszYemini1982}). The matroid structure
itself --- the four matroid axioms --- is delegated to the abstract
$(k, \ell)$-count matroid (\cref{sec:count-matroid}): the rigidity
matroid is the $(k, \ell) = (2, 3)$ specialisation, and
Lov\'asz--Yemini becomes the statement that this combinatorial matroid
coincides with the relation ``$I$ is row-independent at some
placement''. The linear-matroid framing --- building the same matroid
directly from the rigidity-row function at a generic placement via
\texttt{Matroid.ofFun}, with Lov\'asz--Yemini as a matroid equality ---
is the subject of \cref{sec:rigidity-matroid-linear-framing} below.

\begin{definition}[Planar rigidity matroid]
  \label{def:rigidityMatroid}
  \lean{SimpleGraph.rigidityMatroid}
  \leanok
  \uses{def:countMatroid}
  For finite $V$, the \emph{planar rigidity matroid}
  $\mathcal{R}(V)$ is the $(2, 3)$-count matroid on $V$ of
  \cref{def:countMatroid} (which is a matroid since $3 < 4 = 2k$ at
  $k = 2$). Equivalently, its independent sets are exactly the
  $(2, 3)$-sparse edge subsets of $K_{V}$.
\end{definition}

\begin{theorem}[Lov\'asz--Yemini, matroid form]
  \label{thm:rigidityMatroid-indep-iff-rowIndependent}
  \lean{SimpleGraph.rigidityMatroid_indep_iff_edgeSetRowIndependent}
  \leanok
  \uses{def:rigidityMatroid, def:edgeSet-rowIndependent,
    thm:edgeSet-rowIndependent-iff-isSparse}
  Let $V$ be finite and $I \subseteq E(K_{V})$. Then the image of $I$
  in $\mathrm{Sym}_{2}V$ is independent in the planar rigidity
  matroid $\mathcal{R}(V)$ if and only if there exists a placement
  $p : V \to \R^{2}$ at which $I$ is row-independent.
\end{theorem}
\begin{proof}
  \leanok
  \uses{thm:countMatroid-indep-iff,
    thm:edgeSet-rowIndependent-iff-isSparse}
  By \cref{thm:countMatroid-indep-iff}, the image of $I$ is
  independent in $\mathcal{R}(V) = \mathrm{countMatroid}\,V\,2\,3$
  iff it is contained in $E(K_{V})$ (automatic, since $I \subseteq
  E(K_{V})$ already) and $\mathrm{fromEdgeSet}\,I$ is
  $(2, 3)$-sparse. By \cref{thm:edgeSet-rowIndependent-iff-isSparse}
  at $G = K_{V}$, the latter holds iff $I$ is row-independent at some
  placement.
\end{proof}

\subsection{Linear-matroid framing}
\label{sec:rigidity-matroid-linear-framing}

The matroid identification of the previous subsection packages
$\mathcal{R}(V)$ as the $(2, 3)$-count matroid on $\mathrm{Sym}_2 V$ ---
a combinatorial object whose independence is defined purely in terms of
edge-counts on subsets of $V$. An alternative presentation builds the
\emph{same matroid} directly as a \emph{linear} matroid on the
rigidity-row function at a chosen placement, via the
\texttt{Matroid.ofFun} construction of the
\texttt{apnelson1/Matroid} library.

\begin{definition}[Extended rigidity-row function on $\mathrm{Sym}_2 V$]
  \label{def:linearRigidityRow}
  \lean{SimpleGraph.linearRigidityRow}
  \leanok
  \uses{def:rigidityRow}
  Let $p : V \to \R^{d}$ be a placement. Extend
  $\mathrm{rigidityRow}\,K_V\,p : E(K_V) \to \mathrm{Hom}_\R(\mathrm{Framework}\,V\,d, \R)$
  by zero off $E(K_V)$ to obtain
  $\widetilde{\mathrm{row}}_p : \mathrm{Sym}_2 V \to \mathrm{Hom}_\R(\mathrm{Framework}\,V\,d, \R)$.
\end{definition}
\begin{fmlnote}
  The extension is formalised with \texttt{Function.extend} along the
  subtype embedding $E(K_V) \hookrightarrow \mathrm{Sym}_2 V$. Its values
  off $E(K_V)$ do not affect the linear rigidity matroid below, whose
  ground set is $E(K_V)$.
\end{fmlnote}

\begin{definition}[Linear rigidity matroid at a placement]
  \label{def:linearRigidityMatroid}
  \lean{SimpleGraph.linearRigidityMatroid}
  \leanok
  \uses{def:linearRigidityRow}
  Let $p : V \to \R^{d}$ be a placement. The \emph{linear rigidity matroid
  at $p$} is
  $\mathcal{L}_p(V) := \mathtt{Matroid.ofFun}\,\R\,E(K_V)\,\widetilde{\mathrm{row}}_p$,
  the matroid on $\mathrm{Sym}_2 V$ with ground set $E(K_V)$ whose
  independent sets are the linearly independent restrictions of
  $\widetilde{\mathrm{row}}_p$.
\end{definition}

\begin{lemma}[Linear-matroid independence iff row-independence]
  \label{lem:linearRigidityMatroid-indep-iff}
  \lean{SimpleGraph.linearRigidityMatroid_indep_iff_edgeSetRowIndependent}
  \leanok
  \uses{def:linearRigidityMatroid, def:edgeSet-rowIndependent}
  Let $p : V \to \R^{d}$ and $I \subseteq E(K_V)$. Then $I$ is independent
  in $\mathcal{L}_p(V)$ if and only if $I$ is row-independent at $p$ (in
  the sense of \cref{def:edgeSet-rowIndependent} at $G = K_V$).
\end{lemma}
\begin{proof}
  \leanok
  Unfolding the \texttt{Matroid.ofFun} construction
  (\texttt{Matroid.ofFun\_indep\_iff}), the LHS says $I \subseteq E(K_V)$
  (automatic for an $I$ already typed as a subset of $E(K_V)$) together
  with linear independence of $\widetilde{\mathrm{row}}_p|_I$.
  Composing with the injective subtype embedding
  $E(K_V) \hookrightarrow \mathrm{Sym}_2 V$, and noting the agreement
  $\widetilde{\mathrm{row}}_p \circ \mathrm{Subtype.val} =
  \mathrm{rigidityRow}\,K_V\,p$ on $E(K_V)$, rewrites the
  linear-independence condition as
  $\mathrm{LinearIndepOn}\,\R\,(\mathrm{rigidityRow}\,K_V\,p)\,I$;
  the dual-module / function-module bridge
  \texttt{edgeSet\_rowIndependent\_iff\_linearIndepOn\_rigidityRow}
  identifies this with $K_V.\mathrm{EdgeSetRowIndependent}\,p\,I$.
\end{proof}

\begin{lemma}[Uniformly generic placement, dimension $2$]
  \label{lem:exists-uniform-rowIndependent-placement}
  \lean{SimpleGraph.exists_uniform_rowIndependent_placement_dim_two}
  \leanok
  \uses{def:edgeSet-rowIndependent,
    thm:isSparse-exists-rowIndependent-placement}
  Let $V$ be finite. There exists a placement $p : V \to \R^{2}$ at
  which \emph{every} $(2, 3)$-sparse edge subset of $K_V$ is
  simultaneously row-independent: for every $I \subseteq E(K_V)$ with
  $\mathrm{fromEdgeSet}\,I$ $(2, 3)$-sparse, $I$ is row-independent at $p$.
\end{lemma}
\begin{proof}
  \leanok
  Induction on the cardinality of the finite family
  \[
    \mathcal{S} := \{I \subseteq E(K_V) \mid
      \mathrm{fromEdgeSet}\,I\;\text{is}\;(2, 3)\text{-sparse}\}.
  \]
  The base case $\mathcal{S} = \emptyset$ picks any placement. For the
  inductive step, given $p_0$ row-independent on every set in an
  $n$-element subfamily $\mathcal{S}_n \subseteq \mathcal{S}$, and a new
  set $I_0 \in \mathcal{S} \setminus \mathcal{S}_n$,
  \cref{thm:isSparse-exists-rowIndependent-placement} yields a
  placement $q$ row-independent on $I_0$.

  Form the linear-interpolation family
  $p_t := (1 - t)\,p_0 + t\,q = p_0 + t\,(q - p_0)$ for $t \in \R$. The
  rigidity row $\mathrm{row}_{p_t}(e) \in
  \mathrm{Hom}_\R(\mathrm{Framework}\,V\,2, \R)$ is \emph{affine in $t$}
  at each edge $e$ (\texttt{rigidityRow\_add\_smul}: linearity of
  $\mathrm{row}$ in $p$). Hence for each $S \subseteq E(K_V)$, ``$S$ is
  row-independent at $p_t$'' is the non-vanishing of a one-variable
  polynomial in $\R[t]$ (the Gram determinant
  $\det((A + t \cdot B)(A + t \cdot B)^\top)$ of the rigidity-row
  matrices $A, B$ for $p_0$ and $q - p_0$ in a chosen basis of
  $\mathrm{Framework}\,V\,2$). For $S \in \mathcal{S}_n$ this polynomial
  is nonzero at $t = 0$ (the induction hypothesis); for $S = I_0$ it is
  nonzero at $t = 1$.

  Each nonzero polynomial in $\R[t]$ has finitely many roots, so the
  union of bad-$t$ values across $\mathcal{S}_n \cup \{I_0\}$ is finite;
  pick $t$ outside this union and $p_t$ is row-independent on the
  enlarged subfamily. The rectangular linear-independence /
  Gram-determinant equivalence and the cofiniteness-along-an-affine-line
  step are the mirror lemmas
  \texttt{Matrix.linearIndependent\_rows\_iff\_det\_mul\_transpose\_ne\_zero}
  and \texttt{LinearIndependent.finite\_setOf\_not\_along\_affine\_path}.
\end{proof}

\begin{theorem}[Lov\'asz--Yemini, linear-matroid form, dimension $2$]
  \label{thm:linearRigidityMatroid-eq-rigidityMatroid}
  \lean{SimpleGraph.linearRigidityMatroid_eq_rigidityMatroid}
  \leanok
  \uses{def:linearRigidityMatroid, def:rigidityMatroid,
    thm:rigidityMatroid-indep-iff-rowIndependent,
    lem:linearRigidityMatroid-indep-iff,
    lem:exists-uniform-rowIndependent-placement}
  Let $V$ be finite. There exists a placement $p : V \to \R^{2}$ at
  which the linear rigidity matroid coincides with the combinatorial
  planar rigidity matroid: $\mathcal{L}_p(V) = \mathcal{R}(V)$.
\end{theorem}
\begin{proof}
  \leanok
  Take $p$ from
  \cref{lem:exists-uniform-rowIndependent-placement}: every
  $(2, 3)$-sparse edge subset of $K_V$ is row-independent at $p$. The
  two matroids share the ground set
  ($\mathcal{L}_p(V).E = E(K_V) = \mathcal{R}(V).E$), so by
  \texttt{Matroid.ext\_indep} it suffices to identify their
  independent sets. An edge set $I \subseteq E(K_V)$ is independent
  in $\mathcal{L}_p(V)$ iff $I$ is row-independent at $p$
  (\cref{lem:linearRigidityMatroid-indep-iff}); it is independent in
  $\mathcal{R}(V)$ iff $I$ is row-independent at \emph{some} placement
  (\cref{thm:rigidityMatroid-indep-iff-rowIndependent}). The
  $(\Rightarrow)$ direction collapses by taking ``some placement''
  to be $p$; the $(\Leftarrow)$ direction collapses by uniform
  genericity, since every $(2, 3)$-sparse $I$ is row-independent at $p$.
\end{proof}
